\documentclass[11pt,a4paper]{article}
\usepackage[utf8]{utf8}
\usepackage[margin=1in]{geometry}
\usepackage{amsmath,amssymb}
\usepackage{enumitem}

\begin{document}

\begin{center}
    {\Large \textbf{B.Sc. (Hons.) Semester I Examination 2024--25 (NEP)}}\\[4pt]
    {\large \textbf{SKILL ENHANCEMENT COURSE (SEC)}}\\[4pt]
    \textbf{Paper Code: STASE11 : Data Analysis Using Excel}\\[6pt]
    \textbf{Time: 2:30 Hours} \hfill \textbf{Max. Marks: 70}
\end{center}
\hr

\noindent \textbf{Note:} Attempt any FIVE questions. Question No. 1 is compulsory. All questions carry equal marks (14 Marks each).

\begin{enumerate}[label=\textbf{\arabic*.}]
    \item Ages of five children are 2, 5, 8, 3, 4 years and are to be entered in Column A of an Excel sheet using rows 2 to 6 with variable name in Row 1. \hfill [14]
    \begin{enumerate}[label=(\alph*)]
        \item Prepare a layout representation of such an Excel sheet.
        \item Write Excel syntax for: (i) Counting number of observations, (ii) Sum of ages, (iii) Minimum age, (iv) Maximum age, and their returned output values.
    \end{enumerate}

    \item Explain the syntax \texttt{=QUARTILE(array, quart)}. Detail the parameters \texttt{array} and \texttt{quart}. Discuss the shortcut commands \texttt{CTRL+C}, \texttt{CTRL+V}, \texttt{CTRL+SHIFT+J}, \texttt{CTRL+Z}, and \texttt{CTRL+S} in MS Excel. \hfill [14]

    \item Discuss the features and benefits of Pivot Tables in MS Excel. \hfill [14]

    \item Suppose the first 11 natural numbers are entered in cells \texttt{A1:A11}. Evaluate the outputs of the following Excel formulas: \hfill [14]
    \begin{enumerate}[label=(\roman*)]
        \item \texttt{=AVERAGE(A1:A11)}
        \item \texttt{=VAR.P(A1:A11)}
        \item \texttt{=MEDIAN(A1:A11)}
        \item \texttt{=(MAX(A1:A11) - MIN(A1:A11)) / (MAX(A1:A11) + MIN(A1:A11))}
    \end{enumerate}

    \item Explain the four levels of measurement scales (Nominal, Ordinal, Interval, Ratio) with suitable real-world examples. \hfill [14]

    \item The following data represents chemical impurity (ppm) in 25 water samples:
    \[
    11, 19, 24, 30, 12, 20, 25, 29, 15, 21, 24, 31, 16, 23, 25, 26, 32, 17, 22, 26, 35, 18, 24, 18, 27
    \]
    Construct a continuous frequency distribution table with class interval width of 5 and draw a Histogram. \hfill [14]

    \item What are the measures of Central Tendency? Explain methods for calculating Arithmetic Mean for grouped and ungrouped data. \hfill [14]

    \item Based on the 5-point summary below, construct a Box and Whisker Plot. Is the distribution symmetrical/normal? Give reasons.
    \begin{center}
    \begin{tabular}{|c|c|c|c|c|}
    \hline
    \textbf{Minimum} & \textbf{Q1 (First Quartile)} & \textbf{Q2 (Median)} & \textbf{Q3 (Third Quartile)} & \textbf{Maximum} \\ \hline
    5 & 12 & 19 & 26 & 50 \\ \hline
    \end{tabular}
    \end{center} \hfill [14]

    \item Write short notes on any TWO of the following: \hfill [14]
    \begin{enumerate}[label=(\roman*)]
        \item Graphical representation of data in Excel
        \item Uses of Frequency Distribution in data analysis
        \item Pivot Charts and Dashboard creation
        \item Analysis of Variance (ANOVA) in Excel
    \end{enumerate}
\end{enumerate}

\end{document}
